% ch08 — ODE Integration (Runge-Kutta)
\chapter{ODE Integration: Runge-Kutta Methods}
\label{ch:ode-integration}

The ray tracer advances geodesics by integrating a system of ordinary
differential equations. This chapter covers the Runge-Kutta family of
explicit integrators, from fixed-step RK4 through the adaptive
Dormand-Prince method used in the codebase, with attention to the practical
details of step-size control and error estimation.

\section{Explicit Runge-Kutta Methods and Butcher Tableaux}
\label{sec:rk-butcher}

We define the general $s$-stage explicit Runge-Kutta method via its Butcher
tableau and derive the classical RK4 scheme as a special case.

\section{The Dormand-Prince RK45 Pair}
\label{sec:rk-dormand-prince}
\coderef{Geodesic}{Solver}

We present the Dormand-Prince 5(4) embedded pair: a fifth-order solution
propagated at each step, with a fourth-order companion providing the local
error estimate used for step-size control.

\section{First Same As Last (FSAL)}
\label{sec:rk-fsal}

We explain the FSAL optimisation that reuses the last stage of one step as
the first stage of the next, saving one function evaluation per step.

\section{Adaptive Step-Size Control}
\label{sec:rk-adaptive}
\coderef{rk4}{rs}

We describe the PI controller for step-size adaptation, including tolerance
specification, step rejection, and safety factors.

\section{Application: Geodesic Integration}
\label{sec:rk-geodesic-application}

We show how the abstract ODE integrator is applied to the Hamiltonian
geodesic system, connecting the solver parameters to physical accuracy
requirements.

\paragraph{Key References.}
The Butcher-tableau formalism follows \cite{HairerNorsettWanner1993};
Dormand-Prince is from \cite{DormandPrince1980}; adaptive control follows
\cite{HairerWanner1996}.
